% !TeX root = ../main-SMPT_UK.tex
% Проста редакція розділів 4--5; формули та постановка збережені.

\section{Що насправді моделює двомасштабний метод}
\label{sec:analysis}

\subsection{Точність для усередненої мережі}
\label{sec:surrogate}

Якщо визначати час кожної події точно за~\eqref{eq:hazard_integral}, двомасштабний метод сам задає неперервний марковський ланцюг. Він має ті самі локальні правила зараження й видалення, що й повний Gillespie. Відмінність стосується контактів між спільнотами. Замість справжніх ребер метод фактично використовує слабкі зв'язки між усіма парами вузлів двох зв'язаних спільнот. Таку усереднену мережу називаємо сурогатною й позначаємо $\tilde G$.
\par Для сприйнятливого вузла $v\in V_j$ інтенсивність імпорту дорівнює $\sum_{i\neq j}\lambda_{ij}(x)\Pi_{ij}(v\mid x)$. Наступне твердження показує, що це саме інтенсивність зараження вузла в $\tilde G$.

\begin{proposition}[Точна модель для сурогатної мережі]
\label{prop:surrogate}
Нехай $\tilde G$ має ті самі внутрішні контакти та інтенсивності, що й $G$. Для кожної пари спільнот $i\neq j$ замінимо її справжні міжспільнотні канали $E_{ij}$ усіма каналами $V_i\times V_j$. Кожному новому каналу надамо інтенсивність $\tilde\beta_{uv}=\beta_{\mathrm{out}}W_{ij}/(N_iN_j)$. Якщо симулятор вибирає ціль імпорту рівномірно серед сприйнятливих вузлів, ідеальний двомасштабний процес має ті самі інтенсивності всіх переходів, що й повний процес на $\tilde G$. Отже, за однакового початкового стану його розподіл збігається з розподілом повного Gillespie на $\tilde G$.
\end{proposition}

\begin{proof}
Локальні зараження та видалення мають однакові інтенсивності. Для сприйнятливого вузла $v\in V_j$ імовірність вибору ціллю імпорту дорівнює $1/s_j(x)$. Тому його інтенсивність зараження з інших спільнот дорівнює
\[
\begin{aligned}
\sum_{i\neq j}\lambda_{ij}(x)\frac{1}{s_j(x)}
&=\sum_{i\neq j}\beta_{\mathrm{out}}W_{ij}\frac{y_i(x)}{N_i}\frac{s_j(x)}{N_j}\frac{1}{s_j(x)}\\
&=\sum_{i\neq j}\sum_{u\in V_i}\frac{\beta_{\mathrm{out}}W_{ij}}{N_iN_j}\ind\{x_u=I\}.
\end{aligned}
\]
Останній вираз підсумовує інтенсивності всіх каналів $\tilde G$, що ведуть від інфікованих вузлів до $v$. Якщо $s_j(x)=0$, імпорт неможливий в обох моделях. Отже, всі інтенсивності переходів збігаються. \qed
\end{proof}

\begin{corollary}
\label{cor:exact}
Якщо в початковій мережі кожен вузол спільноти $i$ вже з'єднаний із кожним вузлом спільноти $j$, а всі ці канали мають однакову інтенсивність для даної пари спільнот, то $\tilde G=G$. У цьому випадку ідеальний двомасштабний метод точний і для початкової мережі.
\end{corollary}

Загальна сила зв'язку зберігається:
\[
\sum_{(u,v)\in V_i\times V_j}\tilde\beta_{uv}=\beta_{\mathrm{out}}W_{ij}.
\]
Змінюється її розподіл між вузлами (Рисунок~\ref{fig:schematic}c). Тому навіть за точного часу подій результати на $G$ і $\tilde G$ можуть відрізнятися. Реалізація з часовою сіткою додає ще й чисельну похибку. Ці два ефекти можна перевіряти окремо: порівняння з повним Gillespie на $\tilde G$ перевіряє синхронізацію, а порівняння з повним Gillespie на $G$ показує також наслідки заміни контактів.

\subsection{Звідки виникає похибка}
\label{sec:sources}

Заміна справжніх міжспільнотних контактів усередненими змінює три речі: інтенсивність передавання, вузол-одержувач і можливу кількість імпортів. Четверте джерело похибки --- наближене визначення часу подій у реалізації.

\paragraph{1. Не враховано стан справжніх вузлів-мостів.}
У повній мережі інтенсивність передавання з $V_i$ до $V_j$ дорівнює
\begin{equation}
\label{eq:exact_rate}
a^\star_{ij}(x)=\beta_{\mathrm{out}}\sum_{(u,v)\in E_{ij}}w_{uv}\,\ind\{x_u=I,\,x_v=S\},
\end{equation}
Тут у суму входять лише контакти від інфікованого вузла до сприйнятливого. Якщо всі вузли-мости в джерелі ще сприйнятливі, передавання неможливе. У двомасштабному методі інтенсивність залежить від загальної частки інфікованих, тому зараження будь-якого вузла джерела може сприяти експорту. Метод може передати інфекцію, поки жоден справжній міст ще не заражений. Далі називаємо це \emph{обходом мостів} (Розділ~\ref{sec:single_export}).

\paragraph{2. Змінено місце зараження в одержувачі.}
У повній мережі заражається вузол на кінці контакту, через який прийшла інфекція. Двомасштабний метод вибирає будь-який сприйнятливий вузол з однаковою імовірністю. Від цього залежить, як швидко інфекція дійде до вузлів-мостів, що ведуть у наступну спільноту.

\paragraph{3. Можливі додаткові імпорти.}
У SI-моделі вузол після зараження залишається інфікованим. Тому один справжній напрямлений канал може заразити свій цільовий вузол не більше одного разу. Через $|E_{ij}|$ каналів можливі не більше $|E_{ij}|$ заражень; через один канал --- не більше одного.
\par У сурогатній мережі контактів значно більше. Інтенсивність $\lambda_{ij}=\beta_{\mathrm{out}}W_{ij}\ibar_i\sbar_j$ може залишатися додатною після першого імпорту. Очікувана кількість імпортів на інтервалі $[t_1,t_2]$ дорівнює $\mathbb{E}\int_{t_1}^{t_2}\lambda_{ij}(s)\,ds$. Якщо одержувач довго залишається переважно сприйнятливим, метод може додавати зараження, які неможливі через початковий набір контактів.

\paragraph{4. Час події обчислюється наближено.}
Реалізація оцінює інтенсивність за станом у середині інтервалу й використовує цю оцінку для визначення часу імпорту (Розділ~\ref{sec:timing}). Так вона наближено обчислює інтеграл у~\eqref{eq:hazard_integral}. Ця похибка стосується часової сітки й залишається окремою від заміни контактів.

\begin{remark}[Перший імпорт за випадкового розміщення контактів]
\label{rem:placement}
Нехай кінці міжспільнотних каналів вибрано рівномірно й незалежно від внутрішніх графів та інших наборів контактів, як у наших тестах. Якщо спільнота-одержувач ще повністю сприйнятлива, ціль її першого імпорту також рівномірна в повній мережі --- після усереднення за випадковим розміщенням контактів. Двомасштабний метод має такий самий розподіл першої цілі. Для окремої зафіксованої мережі цілі не мусять збігатися. Однаковий розподіл першої цілі також не означає однакового розподілу часу першого імпорту. Порівняння першого й наступних переходів допомагає дослідити наслідки подальшого вибору цілей і додаткових імпортів (Розділ~\ref{sec:res_descriptor}).
\end{remark}

\subsection{Одне передавання інфекції назовні}
\label{sec:single_export}

Щоб окремо розглянути передавання з джерела, приєднаємо до зв’язної спільноти $V_0$ з $n$ вузлів один сприйнятливий вузол $\ell$. Такий зовнішній вузол далі називаємо листком. Він з'єднаний із $k$ вузлами-мостами $B\subset V_0$; кожен канал має інтенсивність $\beta_{\mathrm{out}}$. Усередині джерела діє SI-модель. До зараження листка він не впливає на поширення в джерелі.
\par У двомасштабній моделі $W_{0\ell}=k$, а до зараження листка $\sbar_\ell=1$, тому $\lambda_{0\ell}=\beta_{\mathrm{out}}k\ibar_0$. Позначимо час зараження вузла $u$ через $\tau_u$, а тривалість його перебування інфікованим до часу $t$ --- через $A_u(t)=(t-\tau_u)^+$, де $z^+=\max(z,0)$. Час зараження першого справжнього моста дорівнює $T_B=\min_{u\in B}\tau_u$.

\begin{proposition}[Раніше передавання за випадкового розміщення мостів]
\label{prop:ordering}
Нехай $B$ рівномірно вибрано серед усіх $k$-елементних підмножин $V_0$ незалежно від внутрішнього графа, початкового стану та випадкових подій. Позначимо через $T^{\Full}$ та $T^{\TS}$ часи першого зараження листка у повній мережі та в ідеальному двомасштабному процесі. Тоді
\[
\mathbb{P}(T^{\Full}>t)\ge\mathbb{P}(T^{\TS}>t)
\qquad\text{для всіх }t\ge0.
\]
Це стохастичне впорядкування \citep{shaked2007}: імовірність ще не передати інфекцію до будь-якого заданого часу в двомасштабній моделі не більша. Середній час і кожен квантиль у двомасштабній моделі також не перевищують відповідних значень повної моделі. Твердження усереднює випадкове розміщення мостів; воно не означає раніше передавання в кожній парі запусків або для кожного фіксованого набору мостів.
\end{proposition}

\begin{proof}
До зараження листка джерело розвивається самостійно. Візьмемо його SI-траєкторію $X_0$, продовжену після експорту, та незалежний поріг $E\sim\Exp(1)$. Обидва часи можна записати через цю траєкторію:
\[
\begin{aligned}
T^{\Full}&=\inf\Bigl\{t:\beta_{\mathrm{out}}\sum_{u\in B}A_u(t)\ge E\Bigr\},\\
T^{\TS}&=\inf\Bigl\{t:\beta_{\mathrm{out}}\frac{k}{n}\sum_{u\in V_0}A_u(t)\ge E\Bigr\}.
\end{aligned}
\]
У першому виразі внесок дають лише мости, у другому --- всі вузли з вагою $k/n$. Другий вираз випливає з рівності $\int_0^t\ibar_0(s)\,ds=n^{-1}\sum_u A_u(t)$. За заданих траєкторії $X_0$ та набору мостів $B$ імовірність ще не заразити листок у повній моделі дорівнює $\exp(-\beta_{\mathrm{out}}\sum_{u\in B}A_u(t))$, а в двомасштабній моделі за заданої $X_0$ --- $\exp(-\beta_{\mathrm{out}}(k/n)\sum_uA_u(t))$.
\par За рівномірного незалежного вибору $B$ маємо
\[
\mathbb{E}\Bigl[\sum_{u\in B}A_u(t)\Bigm|X_0\Bigr]
=\frac{k}{n}\sum_u A_u(t).
\]
Функція $z\mapsto e^{-\beta_{\mathrm{out}}z}$ опукла. За нерівністю Єнсена усереднення першого виразу за $B$ дає не менший результат, ніж підстановка усередненого внеску мостів. Отже, $\mathbb{P}(T^{\Full}>t\mid X_0)\ge\mathbb{P}(T^{\TS}>t\mid X_0)$. Усереднення за траєкторіями $X_0$ завершує доведення. \qed
\end{proof}

Різниця є строгою, якщо в момент $t$ тривалості перебування вузлів інфікованими відрізняються з додатною імовірністю й $0<k<n$. Впорядкування зберігається, якщо мости рівномірно вибирати серед вузлів, крім початково зараженого: вилучення вузла, який заражений найдовше, лише зменшує середній внесок мостів. Спільний поріг у доведенні --- спосіб математичного порівняння; в експериментах два симулятори мають незалежні потоки подій.
\par Розкид внеску випадково вибраних мостів дорівнює
\begin{equation}
\label{eq:hazard_variance}
\operatorname{Var}\Bigl(\beta_{\mathrm{out}}\sum_{u\in B}A_u(t)\Bigm| X_0\Bigr)=\beta_{\mathrm{out}}^2\,k\,\frac{n-k}{n-1}\,\operatorname{Var}_u\bigl(A_u(t)\bigr),
\end{equation}
Тут $\operatorname{Var}_u$ --- дисперсія за всіма вузлами джерела, а випадковість ліворуч стосується вибору $B$ за заданої траєкторії $X_0$. Після зараження всіх вузлів $\operatorname{Var}_u(A_u(t))=\operatorname{Var}_u(\tau_u)$. Формула пояснює, чому для точності важливий розкид часів зараження всередині спільноти порівняно з часом $1/\beta_{\mathrm{out}}$. Коли вузли заражаються майже одночасно, усереднення втрачає менше інформації про стан мостів. Саме по собі це ще не дає загальної межі похибки.

Той самий опис дає ймовірність обходу мостів. У повній мережі листок не може заразитися раніше за $T_B$. В ідеальній двомасштабній моделі
\begin{equation}
\label{eq:bypass}
\mathbb{P}\bigl(T^{\TS}<T_B\bigr)=\mathbb{E}\Bigl[1-\exp\Bigl(-\beta_{\mathrm{out}}\frac{k}{n}\int_0^{T_B}y_0(s)\,ds\Bigr)\Bigr],
\end{equation}
Вираз вимірює шанс того, що внеску інших заражених вузлів уже вистачить для експорту, поки всі справжні мости ще сприйнятливі. Що більша накопичена інтенсивність до $T_B$, то більша ця умовна імовірність. Пізнє зараження мостів може її збільшувати. Для джерела-кліки середній час $\mathbb{E}[T_B]$ наведено в Додатку~\ref{app:clique}.

\subsection{Простий показник співвідношення часів}
\label{sec:descriptor}

Формула~\eqref{eq:hazard_variance} пов'язує втрату інформації про мости з розкидом часів зараження вузлів. Додаткові імпорти також залежать від того, як довго одержувач залишається сприйнятливим: їхня кількість визначається накопиченою інтенсивністю передавання. Тому для пояснення результатів порівнюємо два часові масштаби.
\par Нехай $\sigma_T$ --- стандартне відхилення часів зараження вузлів у спільноті. Воно характеризує розкид цих часів, а не весь час зараження спільноти. Середній час передавання через один активний міжспільнотний канал дорівнює $1/\beta_{\mathrm{out}}$. Їхнє відношення дає безрозмірний показник
\begin{equation}
\label{eq:rho}
\rho=\beta_{\mathrm{out}}\,\sigma_T,
\end{equation}
Мале $\rho$ означає, що розкид внутрішніх часів малий порівняно з очікуванням передавання через один канал. Ми перевіряємо, чи допомагає цей показник пояснити похибку; не виводимо з нього загального закону або гарантії точності.
\par Кількість міжспільнотних каналів $k$ теж важлива. У~\eqref{eq:hazard_variance} стандартне відхилення внеску мостів залежить від $\sqrt{k(n-k)/(n-1)}$, а інтенсивність імпортів пропорційна $k$. Тому простого спільного множника для цих механізмів немає. Ми залишаємо $k$ поза означенням $\rho$ та оцінюємо це рішення за даними. У ланцюгових тестах $\sigma_T$ обчислено за часами зараження у \Full{} й усереднено за спільнотами-джерелами та запусками (Розділ~\ref{sec:res_descriptor}). Отже, тут це показник для аналізу вже отриманих результатів; оцінювання лише за локальними симуляторами ще не перевірено.

\section{Як ми перевіряємо метод}
\label{sec:design}

\subsection{Питання і тестові мережі}

Тести побудовано від одного передавання до ланцюга спільнот. Це дає змогу спочатку побачити наслідки усереднення мостів, а потім перевірити, як відмінності змінюються після кількох переходів.
\par \textbf{Кліка--листок.} До джерела зі 100 вузлів, кожен із яких контактує з усіма іншими, приєднано один зовнішній вузол через $k=2$ мости. Це простий тест часу експорту й обходу мостів.
\par \textbf{Випадкова спільнота--листок.} Джерело має $n=1000$ вузлів. Змінюємо його внутрішню структуру, щоб дослідити зв'язок між внутрішнім поширенням і передаванням назовні.
\par \textbf{Ланцюги спільнот.} З'єднуємо $L=10$ спільнот у лінію, по $k$ ребер між сусідніми спільнотами. Інфекція починається у спільноті 0. Ланцюг клік показує випадок швидкого внутрішнього поширення, а перебір розмірів і структур --- зміну та накопичення похибки. Параметри наведено в Таблиці~\ref{tab:design}.
\par Усі тести мають SI-динаміку, одиничні ваги контактів, $\beta_{\mathrm{in}}=\beta_{\mathrm{out}}=\beta$ та крок синхронізації $\tau=1$ (Алгоритм~\ref{alg:scheduler}). Внутрішню структуру описуємо середньою довжиною найкоротшого шляху (ASPL): середньою кількістю ребер у найкоротших шляхах між парами вузлів. ASPL 1 відповідає повному графу.
\par Випадковий зв'язний граф отримуємо з рівномірно випадкового дерева з міченими вузлами, додаючи випадкові ребра. Їхню кількість підбираємо так, щоб ASPL відрізнялася від цілі не більше ніж на 0.1. У цій побудові довші шляхи супроводжуються меншою середньою кількістю сусідів. Наприклад, за цільової ASPL 6 середній степінь становить від 2.08 для $n=50$ до 3.4 для $n=1000$. Тому ці тести не відокремлюють вплив довжини шляхів від щільності контактів.
\par Кінці міжспільнотних ребер вибираємо рівномірно й незалежно для кожного інтерфейсу та повтору.\footnote{У тестах «джерело--листок» мости вибирають серед усіх вузлів, крім одного з фіксованою міткою. Випадкові графи та рівномірно вибраний початковий вузол мають однаковий закон за перемічування вузлів. Тому після такого усереднення виключення фіксованої мітки не змінює розподілу часу експорту, і порівняння з Твердження~\ref{prop:ordering} застосовне.} Початковий заражений вузол випадковий і спільний для двох симуляторів; він може бути мостом. У тесті кліка--листок початковий вузол кліки фіксований.

\begin{table}[H]
\centering
\caption{Параметри тестів. Кожна комірка перебору має 40 реалізацій мережі, по 10 повторів контактів між спільнотами й випадкових подій на кожній: 400 пар запусків. ASPL --- цільова середня довжина найкоротшого шляху всередині спільноти; 1 означає повний граф.}
\label{tab:design}
\small
\begin{tabular}{>{\raggedright\arraybackslash}p{0.17\linewidth}>{\raggedright\arraybackslash}p{0.36\linewidth}>{\raggedright\arraybackslash}p{0.06\linewidth}>{\raggedright\arraybackslash}p{0.08\linewidth}>{\raggedright\arraybackslash}p{0.2\linewidth}}
\toprule
Серія & Структура & $k$ & $\beta$ & Запусків на комірку та модель\\
\midrule
Кліка--листок & кліка-джерело $n=100$ $+$ один листок & 2 & 0.005 & 10\,000\\
Випадкове джерело--листок & джерело $n=1000$, ASPL 1--6 із кроком 0.25, $+$ один листок & 1, 5 & 0.003 & 400 (40 $\times$ 10)\\
Ланцюг клік & $L=10$ клік із $n=1000$ & 1 & 0.003 & 400 (40 $\times$ 10)\\
Ланцюги різних розмірів & $L=10$, $n\in\{50,100,250,1000\}$, ASPL 1--6 із кроком 1 & 1, 5 & 0.003 & 400 (40 $\times$ 10)\\
\bottomrule
\end{tabular}
\end{table}

Кожна пара запусків \Full{} і \TS{} має ті самі граф, міжспільнотні контакти та початковий вузол. Потоки випадкових подій у двох моделях незалежні. Отже, «парні запуски» означають спільну постановку, а не однакові послідовності заражень.
\par Основний горизонт спостереження дорівнює 1000 для кліки--листка та 30\,000 для інших тестів. У двох запусках \Full{} остання спільнота не була заражена до 30\,000, тому їх продовжили до 40\,000 (Додаток~\ref{app:repro}). Усі часи, використані для метрик, були спостережені; замість відсутніх подій не підставляли час завершення запуску. Ці метрики порівнюють часи подій, а не ймовірності завершення на різних горизонтах.
\par Ланцюг клік із $n=1000$ і $k=1$ --- це комірка ASPL 1 у переборі розмірів. Його показано окремо як зручний базовий випадок, але він не є додатковою незалежною серією запусків.

\subsection{Що порівнюємо}

У тестах із листком вимірюємо час його першого зараження $T_{\mathrm{export}}$ та час зараження першого справжнього моста $T_B$. Для ланцюга $T_c$ --- час першого зараження спільноти $c$ від початку запуску. Час досягнення останньої, десятої спільноти дорівнює $T_9$.
\par Три показники відповідають на різні питання:
\begin{itemize}
\item \textbf{Чи прибуває інфекція раніше?} Відносний зсув медіани $b_c=100\,[\mathrm{med}(T^{\TS}_c)-\mathrm{med}(T^{\Full}_c)]/\mathrm{med}(T^{\Full}_c)$ показує напрямок відмінності. Від'ємне значення означає раніше прибуття в \TS{}.
\item \textbf{Наскільки відрізняються розподіли?} Для кожної моделі впорядковуємо $R$ часів за зростанням і обчислюємо $W_1=R^{-1}\sum_r|T^{\TS}_{(r)}-T^{\Full}_{(r)}|$. Це відстань Вассерштайна-1 між емпіричними розподілами. Подаємо її у відсотках від середнього часу \Full{}: $\eW=100\,W_1/\overline{T}^{\Full}$. Тут зіставляються однакові місця у впорядкованих вибірках, а не початкові пари запусків.
\item \textbf{Чи можливий перехід через справжні контакти?} Частка обходу мостів --- відсоток переходів \TS{}, які відбулися до зараження будь-якого справжнього вузла-моста в джерелі.
\end{itemize}
Додатна відстань $\eW$ сама по собі ще не означає систематичної похибки: скінченні вибірки з одного розподілу теж відрізняються.
\par Для локального переходу порівнюємо середні затримки $T_c-T_{c-1}$. Окремо зіставляємо перший перехід з переходами 2--9, для яких сумарна затримка дорівнює $T_9-T_1$. Це допомагає відрізнити похибку окремого переходу від відмінності, накопиченої від початку запуску.

\subsection{Як оцінюємо випадкові відмінності}

Запуски на одному графі не є окремими незалежними реалізаціями мережі. Тому довірчі інтервали обчислюємо блоковим бутстрепом: випадково вибираємо з поверненням цілі мережеві блоки й залишаємо разом усі десять повторів кожного вибраного блока. Для підсумкових показників ланцюга робимо 2000 таких вибірок. Межами 95\% інтервалу є 2.5-й і 97.5-й процентилі отриманих значень.
\par Для $\eW$ також використовуємо перестановочний тест. У кожній парі запусків випадково міняємо мітки \Full{} та \TS{} з імовірністю $1/2$ і знову обчислюємо відстань. Повторюємо це 2000 разів. Якщо за відсутності відмінності між методами мітки в парі можна обмінювати, такий розподіл показує величини відстані, сумісні з випадковими коливаннями скінченної вибірки.
\par Його 95-й процентиль використовуємо як верхній орієнтир випадкової розбіжності. Разом із ним подаємо $p$-значення: воно характеризує, наскільки незвичною є спостережувана відстань за цього припущення. Значення нижче порога не доводить, що моделі однакові. Перевірки окремих конфігурацій не мають поправки на множинні порівняння, тому їх не трактуємо як один спільний тест усієї серії. Коли дані окремих запусків недоступні, наближену очікувану відстань між двома вибірками можна оцінити за формулою з Додатка~\ref{app:floor}; це не заміна перестановочного $p$-значення.

Початкові значення генераторів подій залежать від цільової ASPL, мережевого блока та номера повтору, але не від $n$ чи $k$. Тому різні розміри спільнот і кількості міжспільнотних ребер використовують спільні початкові значення генераторів. Це пов'язує випадкові коливання між конфігураціями; узгоджена тенденція в кількох комірках не є набором незалежних підтверджень.

\subsection{Перевірки повного Gillespie}

Для кліки з рівномірно розміщеними мостами середній час зараження першого моста відомий аналітично (Додаток~\ref{app:clique}). Порівнюємо з ним результати \Full{} для чотирьох розмірів клік і двох кількостей міжспільнотних ребер. У цих перевірках відхилення середнього $T_B$ від формули не перевищувало 2.3 стандартної похибки.
\par За одного ребра після зараження моста листок очікує передавання в середньому $1/\beta=333.3$ одиниці часу. Виміряні середні 326.5--354.0 узгоджувалися з цим значенням у межах випадкової похибки (Таблиця~\ref{tab:verification}). Це перевірка узгодженості реалізації \Full{} з відомими формулами, а не доведення точності двомасштабного наближення.
